%
% Documento: Resumo (Inglês)
\begin{RESUMO}
\thispagestyle{empty}
A rotina de um escritório de advocacia envolve o acompanhamento de processos, o controle de prazos, a organização de documentos e o cadastro de clientes. Quando essas atividades são realizadas de forma manual, 
podem ocorrer perda de prazos, dificuldade no acesso às informações e aumento do tempo gasto em tarefas administrativas. Este trabalho apresenta o desenvolvimento do JurisFlow, uma plataforma web voltada 
para a gestão de escritórios de advocacia, com o objetivo de centralizar informações de clientes, pastas, casos, processos, diligências, prazos e indicadores gerenciais em um único ambiente. 
A metodologia adotada foi estruturada em etapas, incluindo levantamento de requisitos, modelagem do banco de dados, desenvolvimento incremental e verificação das funcionalidades. Para o desenvolvimento, 
foram utilizadas as tecnologias JavaScript, Node.js, PostgreSQL e Docker, além das ferramentas PlantUML, Trello, Git e GitHub para modelagem, organização e controle de versão. Como resultado, 
foi implementado um módulo inicial que contempla o cadastro de usuários e perfis, o gerenciamento de clientes, pastas e casos, o controle de diligências e prazos, notificações automáticas e um 
painel gerencial para acompanhamento das atividades. Os testes realizados permitiram verificar o funcionamento das rotas da API e a persistência dos dados no banco. O principal diferencial do JurisFlow está 
em estruturar as demandas de acordo com a fase de cada caso, diferenciando atividades pré-processuais daquelas que já possuem processo judicial formalizado. Dessa forma, o sistema permite acompanhar e registrar 
as etapas realizadas antes mesmo do protocolo judicial, contribuindo para evitar que os casos dos clientes fiquem paralisados injustificadamente. Conclui-se que o JurisFlow pode contribuir para a organização das 
informações e o acompanhamento das demandas jurídicas, oferecendo maior controle sobre o ciclo completo de cada caso.
\vspace*{0.5cm}\noindent\textbf{Palavras-chave}: Gestão jurídica. Plataforma web. Escritório de advocacia. Engenharia de Software.
\end{RESUMO}

\begin{ABSTRACT}
\thispagestyle{empty}

\vspace*{0.5cm}\noindent\textbf{Keywords}: Legal management. Web platform. Law firm. Software Engineering.
\end{ABSTRACT}
